%TC:ignore
%=======================================
% Shared preamble for every personal statement.
% Each program file sets \school, \pstitle and (optionally) \wordmin and
% \wordlimit BEFORE \input-ing this file; anything left unset falls back to
% the defaults below.
%=======================================
\usepackage{lmodern}          % AER font
\usepackage[T1]{fontenc}      % Better font encoding
\usepackage{amsmath, amssymb}
\usepackage[margin=1in]{geometry}
\usepackage{graphicx}
\usepackage{url}
\usepackage{xcolor}
\usepackage[round]{natbib}
\usepackage{bibentry}         % lets \nobibliography resolve cites without printing a list
\usepackage{hyperref}
\usepackage{fancyhdr}

%=======================================
% Format helpers
%=======================================
\providecommand{\school}{[University name]}
\providecommand{\program}{Economics}
\providecommand{\pstitle}{Personal Statement}
\providecommand{\schoolheader}{\school}   % header line; defaults to \school
% Path from the compiling file's folder back to the repo root (school files
% sit in schools/<school>/tex/; Master_PS overrides this).
\providecommand{\reporoot}{../../../}

% Run-in section heading: bold label ending in a period, body text follows
% on the same line (matches the reference format).
\newcommand{\shead}[1]{\par\vspace{0.85em}\noindent\textbf{#1.}\ }

% Suppress the indent of whichever paragraph starts next. Unlike \noindent
% this survives blank lines/comments in between, and does nothing if no
% paragraph follows (so a trailing \psbreak can't create an empty page).
\newcommand{\psnoindent}{\everypar{{\setbox0\lastbox}\everypar{}}}

% Break between master pieces: a blank line, next line not indented.
\newcommand{\psbreak}{\par\vspace{\baselineskip}\psnoindent}

% Red bracketed note for text still to be written. Each program file tells
% texcount to skip it (see the %TC:macro line in each program file).
\newcommand{\placeholder}[1]{\textcolor{red}{[#1]}}

% Indented paragraphs, no extra vertical space between them.
\setlength{\parindent}{1.5em}
\setlength{\parskip}{0pt}

%=======================================
% Footer: page number centred, word count on the right.
% The count is read from <jobname>.wc, which the texcount build recipe
% regenerates on each build. Text between %TC:ignore and %TC:endignore is
% excluded from the count. The count turns red when outside the limits.
%=======================================
\newread\wcfile
\newif\ifwcknown
\IfFileExists{\jobname.wc}{%
  \openin\wcfile=\jobname.wc
  \read\wcfile to\wcnum
  \closein\wcfile
  \wcknowntrue}{\def\wcnum{??}}

\newif\ifwcbad
\newcommand{\wordcountdisplay}{%
  \wcbadfalse
  \ifwcknown
    \ifdefined\wordlimit\ifnum\wcnum>\wordlimit\relax\wcbadtrue\fi\fi
    \ifdefined\wordmin\ifnum\wcnum<\wordmin\relax\wcbadtrue\fi\fi
  \fi
  \ifwcbad\color{red}\fi
  Word count: \wcnum%
  \ifdefined\wordlimit
  %
  \fi}

\pagestyle{fancy}
\fancyhf{}
\renewcommand{\headrulewidth}{0pt}
\fancyfoot[C]{\thepage}
\fancyfoot[R]{\small\wordcountdisplay}
%TC:endignore
